\subsection{Gesamtarchitektur}
\label{subsec:gesamtarchitektur}

Die untersuchte Baseline ist ein Master-Repository mit gemeinsam gepflegtem
Kern und deklarativen Varianten, kein zur Laufzeit zusammenhängender Monolith.
Profile wählen wiederverwendbare Features aus; optionale Capabilities ergänzen
nur kompatible Profile. Dieses Vorgehen überträgt das Prinzip beherrschter
Variabilität aus Software-Produktlinien auf die Projektstruktur und adressiert
damit die in Kapitel~\ref{sec:anforderungen} geforderte zentrale Wartbarkeit
ohne getrennte Template-Kopien
\parencite{clements2001software-product-lines,kleiveist2026profiles}.

\begin{figure}[H]
  \centering
  \resizebox{\textwidth}{!}{%
  \begin{tikzpicture}[node distance=9mm and 12mm]
    \node[casePrimary,minimum width=32mm] (browser) {Browser};
    \node[casePrimary,minimum width=32mm,right=of browser] (desktop) {Desktop / Tauri};
    \node[caseBox,minimum width=76mm,below=of $(browser)!0.5!(desktop)$] (frontend)
      {gemeinsames Frontend\\Vite + TypeScript};
    \node[caseBox,minimum width=38mm,below=of frontend] (backend) {FastAPI-Backend};
    \node[caseOptional,minimum width=38mm,below=of backend] (database)
      {PostgreSQL\\optionale Persistenz};
    \node[caseBox,minimum width=31mm,left=22mm of backend] (contracts) {Shared Contracts};
    \node[caseBox,minimum width=31mm,right=22mm of backend] (tooling) {Projekt-Tooling};

    \draw[caseFlow] (browser) -- (frontend);
    \draw[caseFlow] (desktop) -- (frontend);
    \draw[caseFlow] (frontend) -- node[right,font=\sffamily\tiny] {HTTP} (backend);
    \draw[caseSecondaryFlow] (backend) -- (database);
    \draw[caseSecondaryFlow,<->] (contracts) -- (backend);
    \draw[caseSecondaryFlow,<->] (contracts) |- (frontend);
    \draw[caseSecondaryFlow] (tooling) -- (backend);
    \draw[caseSecondaryFlow] (tooling) |- (frontend);
  \end{tikzpicture}}
  \caption{Gesamtarchitektur des Technologie-Templates}
  \label{fig:system-architecture}
\end{figure}


Abbildung~\ref{fig:system-architecture} trennt deshalb die vollständige
Baseline von einem abgeleiteten Projekt. Der Generator löst Profil und
Capabilities auf, übernimmt nur die zugehörigen Pfade und hält das Ergebnis in
\texttt{project-profile.toml} fest. Frontend, Backend, Tauri und Persistenz
bleiben eigenständige, teilweise optionale Bausteine
\parencite{kleiveist2026architecture}.

\begin{figure}[H]
  \centering
  \resizebox{\textwidth}{!}{%
  \begin{tikzpicture}[node distance=7mm and 8mm]
    \node[caseBox,minimum width=38mm] (frontend) {\textbf{frontend/}\\Benutzeroberfläche und Client-Build};
    \node[caseBox,minimum width=38mm,right=of frontend] (backend) {\textbf{backend/}\\API und serverseitige Dienste};
    \node[caseBox,minimum width=38mm,right=of backend] (tauri) {\textbf{src-tauri/}\\Desktop-Schicht};
    \node[caseBox,minimum width=38mm,below=of frontend] (shared) {\textbf{shared/}\\gemeinsame Verträge};
    \node[caseBox,minimum width=38mm,right=of shared] (profiles) {\textbf{profiles/}\\Profil-Presets};
    \node[caseBox,minimum width=38mm,right=of profiles] (tools) {\textbf{tools/}\\Projektsteuerung};
    \node[caseBox,minimum width=38mm,below=of shared] (deployment) {\textbf{deployment/}\\Betriebsartefakte};
    \node[caseBox,minimum width=38mm,right=of deployment] (docs) {\textbf{docs/}\\Dokumentation};
    \node[casePrimary,minimum width=38mm,right=of docs] (repository) {Master Repository};

    \draw[caseSecondaryFlow,<->] (frontend) -- (shared);
    \draw[caseSecondaryFlow,<->] (backend) -- (shared);
    \draw[caseSecondaryFlow,<->] (tauri) -- (frontend);
    \draw[caseSecondaryFlow] (profiles) -- (frontend);
    \draw[caseSecondaryFlow] (profiles) -- (backend);
    \draw[caseSecondaryFlow] (profiles) -- (tauri);
    \draw[caseSecondaryFlow] (tools) -- (profiles);
    \draw[caseSecondaryFlow] (deployment) -- (backend);
    \draw[caseSecondaryFlow] (docs) -- (repository);
  \end{tikzpicture}}
  \caption{Komponenten und Verantwortungsgrenzen im Master Repository}
  \label{fig:component-architecture}
\end{figure}


Abbildung~\ref{fig:component-architecture} ordnet die Verantwortungen
konkreten Verzeichnissen zu. \path{frontend/} enthält Darstellung und
HTTP-Client, \path{backend/} die serverseitige API,
\path{src-tauri/} die native Grenze. \path{shared/} hält
frameworkneutrale Artefakte; \path{profiles/} und \path{config/} beschreiben
Struktur beziehungsweise Laufzeitwerte. Tooling und Deployment steuern diese
Komponenten, gehören aber nicht zu deren Produktlaufzeit.

\begin{figure}[H]
  \centering
  \resizebox{0.94\textwidth}{!}{%
  \begin{tikzpicture}[node distance=8mm and 16mm]
    \node[casePrimary,minimum width=32mm] (browser) {Browser};
    \node[casePrimary,minimum width=32mm,right=42mm of browser] (desktop) {Tauri Desktop};
    \node[caseBox,minimum width=32mm,below=of browser] (webfront) {gemeinsames Frontend};
    \node[caseBox,minimum width=32mm,below=of desktop] (deskfront) {gleiches Frontend};
    \node[caseBox,minimum width=39mm,below=15mm of $(webfront)!0.5!(deskfront)$] (api)
      {FastAPI};
    \node[caseOptional,minimum width=39mm,below=of api] (postgres) {PostgreSQL optional};

    \draw[caseFlow] (browser) -- (webfront);
    \draw[caseFlow] (desktop) -- (deskfront);
    \draw[caseFlow] (webfront) -- node[left,font=\sffamily\tiny] {HTTP} (api);
    \draw[caseFlow] (deskfront) -- node[right,font=\sffamily\tiny] {HTTP} (api);
    \draw[caseSecondaryFlow] (api) -- (postgres);

    \begin{scope}[on background layer]
      \node[caseGroup,fit=(browser)(desktop)(postgres),label={[font=\sffamily\scriptsize,text=black!60]below:Komponentenbestand abhängig vom gewählten Profil}] {};
    \end{scope}
  \end{tikzpicture}}
  \caption{Profilabhängige Laufzeitarchitektur}
  \label{fig:runtime-architecture}
\end{figure}


Zur Laufzeit wird derselbe Vite-Build entweder vom Browser oder in einer
Tauri-Webview geladen (Abbildung~\ref{fig:runtime-architecture}). Nur
backendfähige Profile ergänzen die Kommunikation über die HTTP-Grenze zu
FastAPI; ausschließlich das Backend darf eine aktivierte Persistenzschicht
ansprechen. Damit bleiben Client, Server und Datenhaltung getrennt test- und
erweiterbar, ohne dass jedes Profil alle Laufzeiteinheiten enthalten muss.
